\documentclass[../../../uem_tok_q.tex]{subfiles}

\begin{document}
    白玉2個が横に並んでいる．投げたとき表と裏のでる確率がそれぞれ $\myfraca{1}{2}$ %
    のコインを用いて，次の手順($\ast$)をくり返し，白玉または黒玉を横一列に並べていく．

    \begin{itemize}[label={手順($\ast$)}, leftmargin=5\zw, labelsep=.7\zw, topsep=4pt]
        \item コインを投げ，表が出たら白玉，裏がでたら黒玉を，%
            それまでに並べられている一番右にある玉の右隣におく．%
            そして，新しくおいた玉の色がその1つ左の玉の色と異なり，%
            かつ2つ左の玉の色と一致するときには，%
            新しくおいた玉の1つ左の玉を新しくおいた玉と同じ色の玉にとりかえる．
    \end{itemize}
    例えば，手順($\ast$)を2回行いコインが裏，表の順にでた場合には，白玉が4つ並ぶ．%
    正の整数$n$に対して，手順($\ast$)を$n$回行った時点での$(n+2)$個の玉の並び方を考える．

    \begin{enumerate}
        \item $n=3$ のとき，右から2番目の玉が白玉である確率を求めよ．
        \item $n$を正の整数とする．右から2番目の玉が白玉である確率を求めよ．
        \item $n$を正の整数とする．右から1番目と2番目の玉がともに白玉である確率を求めよ．
    \end{enumerate}
   
        \setlength{\lineskip}{0pt}
        \setlength{\lineskiplimit}{0pt}

    $n$を正の整数とする．右から1番目と2番目の玉がともに白玉である確率を求めよ．

          
\end{document}
